\begin{center}
\LARGE
A Mathematica Package for Cooperative Problem Solving in
                   Computational Mechanics

\end{center}

\begin{center}
\Large
J. Korelc

\end{center}

\noindent
Date:  July 20th (Saturday)\\
Time:  09:30-10:00\\

\begin{center}
\large
Abstract
\end{center}

he current state of symbolic support in an advanced mechanical analysis is still very low. Using
the general symbolic systems at high abstract level leads to extensive expression growth and
unusable results. When the low level knowledge about solution is established and the sequence of
matrix and vector operations, which lead to the desired elements characteristics is known, then the
symbolic approach is not needed any more. Investigations and experiments with different methods
like symbolic and algebraic systems, automatic differentiation tools, automatic theorem proving
tools and problem solving environments, show that efficient treatment of complex problems can
be obtained only by a cooperation of all these techniques. 

A new Mathematica package, named SMS (Symbolic Mechanics System), which represents a
pragmatic approach to the problem of combining different techniques, will be introduced. A new
approach is called Simultaneous Stochastic Simplification of numerical code. It combines a
general computer algebra system Mathematica with an automatic differentiation technique and
theorem proving by examples in order to automatically generate the finite element code. To
alleviate the problem of growth of expressions and redundant calculations, simultaneous
simplification of symbolic expressions is used. It is based on a stochastic evaluation of the
formulas instead on a conventional pattern matching technique. 

The new package was system has already been successfully used for the development of the new
2D and 3D elements based on a modified enhanced strain method. Among the other problems, the
well known hour-glassing problem of enhanced strain elements in a presence of large
deformations was successfully solved with the help of the new system. It was also applied to
develop a `spline' finite strip element for nonlinear analysis of the prismatic shell structures with
the non uniform cross section. 

